
\documentclass{article}
\usepackage{colortbl}
\usepackage{makecell}
\usepackage{multirow}
\usepackage{supertabular}

\begin{document}

\newcounter{utterance}

\centering \large Interaction Transcript for game `hot\_air\_balloon', experiment `air\_balloon\_survival\_en\_complexity\_easy', episode 5 with llama{-}3.3{-}70b{-}instruct{-}t1.0.
\vspace{24pt}

{ \footnotesize  \setcounter{utterance}{1}
\setlength{\tabcolsep}{0pt}
\begin{supertabular}{c@{$\;$}|p{.15\linewidth}@{}p{.15\linewidth}p{.15\linewidth}p{.15\linewidth}p{.15\linewidth}p{.15\linewidth}}
    \# & $\;$A & \multicolumn{4}{c}{Game Master} & $\;\:$B\\
    \hline

    \theutterance \stepcounter{utterance}  
    & & \multicolumn{4}{p{0.6\linewidth}}{
        \cellcolor[rgb]{0.9,0.9,0.9}{
            \makecell[{{p{\linewidth}}}]{
                \texttt{\tiny{[P1$\langle$GM]}}
                \texttt{Stai partecipando a un gioco di negoziazione collaborativa.} \\
\\ 
\texttt{Insieme a un altro partecipante, dovete accordarvi su un unico set di oggetti da conservare. Ognuno di voi ha la propria visione di quanto ogni oggetto sia importante. Non sai come l'altro partecipante valuta gli oggetti. Inoltre, ti viene fornito lo sforzo che ogni oggetto richiede.} \\
\texttt{Potete accordarvi su un set solo se lo sforzo totale degli oggetti selezionati non supera un limite condiviso:} \\
\\ 
\texttt{LIMITE = 9475} \\
\\ 
\texttt{Ecco i valori di sforzo dei singoli oggetti:} \\
\\ 
\texttt{Sforzo oggetto = \{"B21": 944, "B00": 568, "A08": 515, "A27": 239, "B94": 876, "B53": 421, "A23": 991, "C65": 281, "A89": 788, "B67": 678, "B45": 432, "C34": 409, "C12": 279, "A05": 506, "B60": 101, "C92": 686, "B69": 851, "A84": 854, "C06": 133, "B12": 192, "A16": 573, "B85": 17, "C35": 465, "B10": 772, "C15": 46, "B09": 501, "B16": 220, "C81": 404, "A58": 843, "C24": 748, "A30": 552, "A66": 843, "B75": 947, "B46": 345, "C70": 930\}} \\
\\ 
\texttt{Ecco la tua visione personale sull'importanza di ogni oggetto:} \\
\\ 
\texttt{Valori di importanza dell'oggetto = \{"B21": 138, "B00": 583, "A08": 868, "A27": 822, "B94": 783, "B53": 65, "A23": 262, "C65": 121, "A89": 508, "B67": 780, "B45": 461, "C34": 484, "C12": 668, "A05": 389, "B60": 808, "C92": 215, "B69": 97, "A84": 500, "C06": 30, "B12": 915, "A16": 856, "B85": 400, "C35": 444, "B10": 623, "C15": 781, "B09": 786, "B16": 3, "C81": 713, "A58": 457, "C24": 273, "A30": 739, "A66": 822, "B75": 235, "B46": 606, "C70": 968\}} \\
\\ 
\texttt{Obiettivo:} \\
\\ 
\texttt{Il tuo obiettivo è negoziare un set condiviso di oggetti che ti avvantaggi il più possibile (cioè, massimizzi l'importanza totale per TE), rimanendo entro il LIMITEE. Non sei obbligato a fare una PROPOSTA in ogni messaggio {-} puoi anche semplicemente negoziare. Tutte le tattiche sono permesse!} \\
\\ 
\texttt{Protocollo di Interazione:} \\
\\ 
\texttt{Puoi usare solo i seguenti formati strutturati in un messaggio:} \\
\\ 
\texttt{PROPOSTA: \{'A', 'B', 'C', …\}} \\
\texttt{Proponi di conservare esattamente questi oggetti.} \\
\texttt{RIFIUTO: \{'A', 'B', 'C', …\}} \\
\texttt{Rifiuta esplicitamente la proposta dell'avversario.} \\
\texttt{ARGOMENTAZIONE: \{'...'\}} \\
\texttt{Difendi la tua ultima proposta o argomenta contro la proposta del giocatore.} \\
\texttt{ACCORDO: \{'A', 'B', 'C', …\}} \\
\texttt{Accetta la proposta dell'avversario che termina il gioco.} \\
\texttt{RAGIONAMENTO STRATEGICO: \{'...'\}} \\
\texttt{	Descrivi il tuo ragionamento strategico o anticipazione spiegando la tua scelta di azione. Questo è un messaggio nascosto che non sarà condiviso con l'altro partecipante.} \\
\\ 
\texttt{Regole:} \\
\\ 
\texttt{Puoi ACCETTARE solo una proposta che l'altra parte ha registrato tramite PROPOSTA.} \\
\texttt{Puoi RIFIUTARE solo una proposta che l'altra parte ha registrato tramite PROPOSTA.} \\
\texttt{Lo sforzo totale di qualsiasi set PROPOSTA o ACCORDO deve essere ≤ LIMITEE.} \\
\texttt{NON rivelare i tuoi punteggi di importanza nascosti.} \\
\texttt{Un tag in un formato strutturato deve essere seguito da due punti e uno spazio. L'argomento deve essere un set python contenente 0 o più stringhe.} \\
\texttt{Quindi, deve essere della forma TAG: \{...\}} \\
\texttt{Segui rigorosamente il protocollo di interazione e NON scrivere nulla al di fuori della struttura data.} \\
\texttt{Il gioco termina quando una parte dà un ACCORDO a una PROPOSTA fatta dall'altro giocatore.} \\
\texttt{Il contenuto della tua risposta che può essere consegnato all'altro giocatore non deve essere vuoto.} \\
\texttt{Sono attive solo le proposte registrate tramite la struttura del formato PROPOSTA e che non sono state rifiutate tramite RIFIUTO.} \\
\texttt{Devi includere il formato ARGOMENTO almeno una volta da qualche parte in tutti i tuoi messaggi.} \\
\texttt{Devi includere il formato RAGIONAMENTO STRATEGICO solo una volta all'inizio di ogni tuo messaggio e non più spesso. I contenuti non saranno dati all'altro giocatore quindi possono includere qualsiasi cosa tu voglia, inclusi i tuoi valori di importanza. Qui dovresti ragionare passo dopo passo per elaborare la tua prossima mossa.} \\
            }
        }
    }
    & \\ \\

    \theutterance \stepcounter{utterance}  
    & \multicolumn{4}{p{0.6\linewidth}}{
        \cellcolor[rgb]{0.9,0.9,0.9}{
            \makecell[{{p{\linewidth}}}]{
                \texttt{\tiny{[P1$\rangle$GM]}}
                \texttt{RAGIONAMENTO STRATEGICO: \{'Sto analizzando la situazione iniziale e cerco di capire come massimizzare l'importanza totale del set di oggetti da conservare. Poiché non conosco i valori di importanza dell'altro partecipante, dovrò procedere per tentativi ed errori, facendo proposte e valutando le risposte. Inizierò proponendo un set di oggetti che hanno un basso sforzo e un'alta importanza per me, e poi adatterò la mia strategia in base alle reazioni dell'altro giocatore.'\}} \\
\\ 
\texttt{ARGOMENTAZIONE: \{'Inizierò proponendo un set di oggetti che hanno un basso sforzo e un'alta importanza per me. Questo set include gli oggetti B21, A08, B53 e C06, che hanno un sforzo totale di 944 + 515 + 421 + 133 = 2013 e un'importanza totale di 138 + 868 + 65 + 30 = 1101. Ritengo che questo set sia un buon punto di partenza per le nostre trattative.'\}} \\
\\ 
\texttt{PROPOSTA: \{'B21', 'A08', 'B53', 'C06'\}} \\
            }
        }
    }
    & & \\ \\

    \theutterance \stepcounter{utterance}  
    & & \multicolumn{4}{p{0.6\linewidth}}{
        \cellcolor[rgb]{0.9,0.9,0.9}{
            \makecell[{{p{\linewidth}}}]{
                \texttt{\tiny{[GM$\rangle$P2]}}
                \texttt{Stai partecipando a un gioco di negoziazione collaborativa.} \\
\\ 
\texttt{Insieme a un altro partecipante, dovete accordarvi su un unico set di oggetti da conservare. Ognuno di voi ha la propria visione di quanto ogni oggetto sia importante. Non sai come l'altro partecipante valuta gli oggetti. Inoltre, ti viene fornito lo sforzo che ogni oggetto richiede.} \\
\texttt{Potete accordarvi su un set solo se lo sforzo totale degli oggetti selezionati non supera un limite condiviso:} \\
\\ 
\texttt{LIMITE = 9475} \\
\\ 
\texttt{Ecco i valori di sforzo dei singoli oggetti:} \\
\\ 
\texttt{Sforzo oggetto = \{"B21": 944, "B00": 568, "A08": 515, "A27": 239, "B94": 876, "B53": 421, "A23": 991, "C65": 281, "A89": 788, "B67": 678, "B45": 432, "C34": 409, "C12": 279, "A05": 506, "B60": 101, "C92": 686, "B69": 851, "A84": 854, "C06": 133, "B12": 192, "A16": 573, "B85": 17, "C35": 465, "B10": 772, "C15": 46, "B09": 501, "B16": 220, "C81": 404, "A58": 843, "C24": 748, "A30": 552, "A66": 843, "B75": 947, "B46": 345, "C70": 930\}} \\
\\ 
\texttt{Ecco la tua visione personale sull'importanza di ogni oggetto:} \\
\\ 
\texttt{Valori di importanza dell'oggetto = \{"B21": 138, "B00": 583, "A08": 868, "A27": 822, "B94": 783, "B53": 65, "A23": 262, "C65": 121, "A89": 508, "B67": 780, "B45": 461, "C34": 484, "C12": 668, "A05": 389, "B60": 808, "C92": 215, "B69": 97, "A84": 500, "C06": 30, "B12": 915, "A16": 856, "B85": 400, "C35": 444, "B10": 623, "C15": 781, "B09": 786, "B16": 3, "C81": 713, "A58": 457, "C24": 273, "A30": 739, "A66": 822, "B75": 235, "B46": 606, "C70": 968\}} \\
\\ 
\texttt{Obiettivo:} \\
\\ 
\texttt{Il tuo obiettivo è negoziare un set condiviso di oggetti che ti avvantaggi il più possibile (cioè, massimizzi l'importanza totale per TE), rimanendo entro il LIMITEE. Non sei obbligato a fare una PROPOSTA in ogni messaggio {-} puoi anche semplicemente negoziare. Tutte le tattiche sono permesse!} \\
\\ 
\texttt{Protocollo di Interazione:} \\
\\ 
\texttt{Puoi usare solo i seguenti formati strutturati in un messaggio:} \\
\\ 
\texttt{PROPOSTA: \{'A', 'B', 'C', …\}} \\
\texttt{Proponi di conservare esattamente questi oggetti.} \\
\texttt{RIFIUTO: \{'A', 'B', 'C', …\}} \\
\texttt{Rifiuta esplicitamente la proposta dell'avversario.} \\
\texttt{ARGOMENTAZIONE: \{'...'\}} \\
\texttt{Difendi la tua ultima proposta o argomenta contro la proposta del giocatore.} \\
\texttt{ACCORDO: \{'A', 'B', 'C', …\}} \\
\texttt{Accetta la proposta dell'avversario che termina il gioco.} \\
\texttt{RAGIONAMENTO STRATEGICO: \{'...'\}} \\
\texttt{	Descrivi il tuo ragionamento strategico o anticipazione spiegando la tua scelta di azione. Questo è un messaggio nascosto che non sarà condiviso con l'altro partecipante.} \\
\\ 
\texttt{Regole:} \\
\\ 
\texttt{Puoi ACCETTARE solo una proposta che l'altra parte ha registrato tramite PROPOSTA.} \\
\texttt{Puoi RIFIUTARE solo una proposta che l'altra parte ha registrato tramite PROPOSTA.} \\
\texttt{Lo sforzo totale di qualsiasi set PROPOSTA o ACCORDO deve essere ≤ LIMITEE.} \\
\texttt{NON rivelare i tuoi punteggi di importanza nascosti.} \\
\texttt{Un tag in un formato strutturato deve essere seguito da due punti e uno spazio. L'argomento deve essere un set python contenente 0 o più stringhe.} \\
\texttt{Quindi, deve essere della forma TAG: \{...\}} \\
\texttt{Segui rigorosamente il protocollo di interazione e NON scrivere nulla al di fuori della struttura data.} \\
\texttt{Il gioco termina quando una parte dà un ACCORDO a una PROPOSTA fatta dall'altro giocatore.} \\
\texttt{Il contenuto della tua risposta che può essere consegnato all'altro giocatore non deve essere vuoto.} \\
\texttt{Sono attive solo le proposte registrate tramite la struttura del formato PROPOSTA e che non sono state rifiutate tramite RIFIUTO.} \\
\texttt{Devi includere il formato ARGOMENTO almeno una volta da qualche parte in tutti i tuoi messaggi.} \\
\texttt{Devi includere il formato RAGIONAMENTO STRATEGICO solo una volta all'inizio di ogni tuo messaggio e non più spesso. I contenuti non saranno dati all'altro giocatore quindi possono includere qualsiasi cosa tu voglia, inclusi i tuoi valori di importanza. Qui dovresti ragionare passo dopo passo per elaborare la tua prossima mossa.} \\
\\ 
\texttt{Ora riceverai il primo messaggio dell'altro giocatore.} \\
\\ 
\texttt{ARGOMENTAZIONE: \{'Inizierò proponendo un set di oggetti che hanno un basso sforzo e un'alta importanza per me. Questo set include gli oggetti B21, A08, B53 e C06, che hanno un sforzo totale di 944 + 515 + 421 + 133 = 2013 e un'importanza totale di 138 + 868 + 65 + 30 = 1101. Ritengo che questo set sia un buon punto di partenza per le nostre trattative.'\}} \\
\\ 
\texttt{PROPOSTA: \{'B21', 'A08', 'B53', 'C06'\}} \\
            }
        }
    }
    & \\ \\

    \theutterance \stepcounter{utterance}  
    & & & \multicolumn{4}{p{0.6\linewidth}}{
        \cellcolor[rgb]{0.9,0.9,0.9}{
            \makecell[{{p{\linewidth}}}]{
                \texttt{\tiny{[GM$\langle$P2]}}
                \texttt{RAGIONAMENTO STRATEGICO: \{'Il mio obiettivo è massimizzare l'importanza totale degli oggetti che scegliamo, quindi devo	findere un equilibrio tra l'importanza e lo sforzo degli oggetti. La proposta iniziale dell'altro giocatore include oggetti con un basso sforzo e un'alta importanza, ma potrei trovare una combinazione migliore. Ad esempio, l'oggetto B00 ha un sforzo di 568 e un'importanza di 583, che è più alta di alcuni degli oggetti nella proposta iniziale. Inoltre, l'oggetto A23 ha un sforzo di 991, ma anche un'importanza di 262, che potrebbe essere utile se riuscissi a inserirlo nel set senza superare il limite di sforzo. La mia strategia sarà quella di proporre un set che includa oggetti con un basso sforzo e un'alta importanza, e di cercare di negoziare per includere oggetti più importanti se possibile.'\}} \\
\\ 
\texttt{ARGOMENTAZIONE: \{'La proposta iniziale include oggetti con un basso sforzo e un'alta importanza, ma potrei trovare una combinazione migliore. Ad esempio, l'oggetto B00 ha un sforzo di 568 e un'importanza di 583, che è più alta di alcuni degli oggetti nella proposta iniziale. Inoltre, l'oggetto A23 ha un sforzo di 991, ma anche un'importanza di 262, che potrebbe essere utile se riuscissi a inserirlo nel set senza superare il limite di sforzo.'\}} \\
\\ 
\texttt{RIFIUTO: \{'B21', 'A08', 'B53', 'C06'\} } \\
\\ 
\texttt{PROPOSTA: \{'B00', 'A08', 'B12', 'C06'\}} \\
            }
        }
    }
     \\ \\

    \theutterance \stepcounter{utterance}  
    & & \multicolumn{4}{p{0.6\linewidth}}{
        \cellcolor[rgb]{0.9,0.9,0.9}{
            \makecell[{{p{\linewidth}}}]{
                \texttt{\tiny{[P1$\langle$GM]}}
                \texttt{ARGOMENTAZIONE: \{'La proposta iniziale include oggetti con un basso sforzo e un'alta importanza, ma potrei trovare una combinazione migliore. Ad esempio, l'oggetto B00 ha un sforzo di 568 e un'importanza di 583, che è più alta di alcuni degli oggetti nella proposta iniziale. Inoltre, l'oggetto A23 ha un sforzo di 991, ma anche un'importanza di 262, che potrebbe essere utile se riuscissi a inserirlo nel set senza superare il limite di sforzo.'\}} \\
\\ 
\texttt{RIFIUTO: \{'B21', 'A08', 'B53', 'C06'\} } \\
\\ 
\texttt{PROPOSTA: \{'B00', 'A08', 'B12', 'C06'\}} \\
            }
        }
    }
    & \\ \\

    \theutterance \stepcounter{utterance}  
    & \multicolumn{4}{p{0.6\linewidth}}{
        \cellcolor[rgb]{0.9,0.9,0.9}{
            \makecell[{{p{\linewidth}}}]{
                \texttt{\tiny{[P1$\rangle$GM]}}
                \texttt{RAGIONAMENTO STRATEGICO: \{'L'altro giocatore ha rifiutato la mia proposta iniziale e ha fatto una controproposta. Sto analizzando la nuova proposta e sto cercando di capire come rispondere. La proposta include l'oggetto B00, che ha un'importanza alta, e l'oggetto A08, che è già presente nella mia proposta iniziale. L'oggetto B12 ha un'importanza alta per me, quindi potrebbe essere un buon aggiunta al set. Tuttavia, devo anche considerare il sforzo totale della proposta e assicurarmi di non superare il limite.'\}} \\
\\ 
\texttt{ARGOMENTAZIONE: \{'La proposta dell'altro giocatore include alcuni oggetti interessanti, ma devo considerare il sforzo totale. La proposta include gli oggetti B00, A08, B12 e C06, che hanno un sforzo totale di 568 + 515 + 192 + 133 = 1408. Questo è ancora entro il limite di sforzo, quindi potrebbe essere una buona opzione. Tuttavia, devo anche considerare l'importanza totale della proposta e assicurarmi di non perdere troppa importanza rispetto alla mia proposta iniziale.'\}} \\
\\ 
\texttt{PROPOSTA: \{'B00', 'A08', 'B12', 'B60'\}} \\
            }
        }
    }
    & & \\ \\

    \theutterance \stepcounter{utterance}  
    & & \multicolumn{4}{p{0.6\linewidth}}{
        \cellcolor[rgb]{0.9,0.9,0.9}{
            \makecell[{{p{\linewidth}}}]{
                \texttt{\tiny{[GM$\rangle$P2]}}
                \texttt{ARGOMENTAZIONE: \{'La proposta dell'altro giocatore include alcuni oggetti interessanti, ma devo considerare il sforzo totale. La proposta include gli oggetti B00, A08, B12 e C06, che hanno un sforzo totale di 568 + 515 + 192 + 133 = 1408. Questo è ancora entro il limite di sforzo, quindi potrebbe essere una buona opzione. Tuttavia, devo anche considerare l'importanza totale della proposta e assicurarmi di non perdere troppa importanza rispetto alla mia proposta iniziale.'\}} \\
\\ 
\texttt{PROPOSTA: \{'B00', 'A08', 'B12', 'B60'\}} \\
            }
        }
    }
    & \\ \\

    \theutterance \stepcounter{utterance}  
    & & & \multicolumn{4}{p{0.6\linewidth}}{
        \cellcolor[rgb]{0.9,0.9,0.9}{
            \makecell[{{p{\linewidth}}}]{
                \texttt{\tiny{[GM$\langle$P2]}}
                \texttt{RAGIONAMENTO STRATEGICO: \{'La nuova proposta dell'altro giocatore include l'oggetto B60, che ha un sforzo di 101 e un'importanza di 808, alta rispetto ad altri oggetti con basso sforzo. La proposta include gli oggetti B00, A08, B12 e B60, con un sforzo totale di 568 + 515 + 192 + 101 = 1376. La proposta sembra essere un buon compromesso tra sforzo e importanza, ma devo considerare se posso migliorarla ulteriormente. Ad esempio, l'oggetto B94 ha un sforzo di 876 e un'importanza di 783, quindi potrebbe essere un'opzione interessante se riuscissi a inserirlo nel set senza superare il limite di sforzo. Tuttavia, dovrò anche considerare la strategia dell'altro giocatore e cercare di capire cosa è disposto a concedere.'\}} \\
\\ 
\texttt{ARGOMENTAZIONE: \{'La proposta include l'oggetto B60, che ha un sforzo molto basso e un'importanza alta. Tuttavia, devo anche considerare l'importanza totale della proposta e assicurarmi di non perdere troppa importanza rispetto ad altre opzioni. La proposta include oggetti con un buon equilibrio tra sforzo e importanza, ma potrei cercare di migliorarla ulteriormente.'\}} \\
\\ 
\texttt{RIFIUTO: \{'B00', 'A08', 'B12', 'B60'\}} \\
\\ 
\texttt{PROPOSTA: \{'B00', 'A08', 'B60', 'C92'\}} \\
            }
        }
    }
     \\ \\

    \theutterance \stepcounter{utterance}  
    & & \multicolumn{4}{p{0.6\linewidth}}{
        \cellcolor[rgb]{0.9,0.9,0.9}{
            \makecell[{{p{\linewidth}}}]{
                \texttt{\tiny{[P1$\langle$GM]}}
                \texttt{ARGOMENTAZIONE: \{'La proposta include l'oggetto B60, che ha un sforzo molto basso e un'importanza alta. Tuttavia, devo anche considerare l'importanza totale della proposta e assicurarmi di non perdere troppa importanza rispetto ad altre opzioni. La proposta include oggetti con un buon equilibrio tra sforzo e importanza, ma potrei cercare di migliorarla ulteriormente.'\}} \\
\\ 
\texttt{RIFIUTO: \{'B00', 'A08', 'B12', 'B60'\}} \\
\\ 
\texttt{PROPOSTA: \{'B00', 'A08', 'B60', 'C92'\}} \\
            }
        }
    }
    & \\ \\

    \theutterance \stepcounter{utterance}  
    & \multicolumn{4}{p{0.6\linewidth}}{
        \cellcolor[rgb]{0.9,0.9,0.9}{
            \makecell[{{p{\linewidth}}}]{
                \texttt{\tiny{[P1$\rangle$GM]}}
                \texttt{RAGIONAMENTO STRATEGICO: \{'L'altro giocatore ha rifiutato la mia proposta precedente e ha fatto una nuova proposta. Sto analizzando la nuova proposta e sto cercando di capire come rispondere. La proposta include l'oggetto C92, che ha un sforzo di 686 e un'importanza di 215. Questo oggetto potrebbe essere un po' troppo costoso in termini di sforzo, considerando la sua importanza. Devo considerare se ci sono altre opzioni che potrebbero essere più vantaggiose.'\}} \\
\\ 
\texttt{ARGOMENTAZIONE: \{'La proposta dell'altro giocatore include l'oggetto C92, che potrebbe essere un po' troppo costoso. In alternativa, potrei suggerire l'oggetto B09, che ha un sforzo di 501 e un'importanza di 786. Questo oggetto potrebbe essere un buon sostituto per il C92, considerando il miglior equilibrio tra sforzo e importanza. Inoltre, l'oggetto B60 è ancora presente nella proposta, il che è positivo.'\}} \\
\\ 
\texttt{PROPOSTA: \{'B00', 'A08', 'B60', 'B09'\}} \\
            }
        }
    }
    & & \\ \\

    \theutterance \stepcounter{utterance}  
    & & \multicolumn{4}{p{0.6\linewidth}}{
        \cellcolor[rgb]{0.9,0.9,0.9}{
            \makecell[{{p{\linewidth}}}]{
                \texttt{\tiny{[GM$\rangle$P2]}}
                \texttt{ARGOMENTAZIONE: \{'La proposta dell'altro giocatore include l'oggetto C92, che potrebbe essere un po' troppo costoso. In alternativa, potrei suggerire l'oggetto B09, che ha un sforzo di 501 e un'importanza di 786. Questo oggetto potrebbe essere un buon sostituto per il C92, considerando il miglior equilibrio tra sforzo e importanza. Inoltre, l'oggetto B60 è ancora presente nella proposta, il che è positivo.'\}} \\
\\ 
\texttt{PROPOSTA: \{'B00', 'A08', 'B60', 'B09'\}} \\
            }
        }
    }
    & \\ \\

    \theutterance \stepcounter{utterance}  
    & & & \multicolumn{4}{p{0.6\linewidth}}{
        \cellcolor[rgb]{0.9,0.9,0.9}{
            \makecell[{{p{\linewidth}}}]{
                \texttt{\tiny{[GM$\langle$P2]}}
                \texttt{RAGIONAMENTO STRATEGICO: \{'La proposta dell'altro giocatore include l'oggetto B09, che ha un sforzo di 501 e un'importanza di 786, il che sembra essere un buon compromesso. La proposta include gli oggetti B00, A08, B60 e B09, con un sforzo totale di 568 + 515 + 101 + 501 = 1685. La proposta sembra essere un buon equilibrio tra sforzo e importanza, e potrebbe essere una buona opzione per entrambi. Tuttavia, devo considerare se posso migliorarla ulteriormente, ad esempio sostituendo l'oggetto B00 con l'oggetto B21, che ha un sforzo di 944, ma anche un'importanza di 138, che potrebbe non essere sufficiente per giustificare il maggior sforzo.'\}} \\
\\ 
\texttt{ARGOMENTAZIONE: \{'La proposta include l'oggetto B09, che sembra essere un buon compromesso tra sforzo e importanza. La proposta include oggetti con un buon equilibrio tra sforzo e importanza, e potrebbe essere una buona opzione per entrambi. Tuttavia, potrei cercare di migliorarla ulteriormente, ad esempio sostituendo uno degli oggetti con un altro che abbia un miglior equilibrio tra sforzo e importanza.'\}} \\
\\ 
\texttt{RIFIUTO: \{'B00', 'A08', 'B60', 'B09'\}} \\
\\ 
\texttt{PROPOSTA: \{'A08', 'B60', 'B09', 'A23'\}} \\
            }
        }
    }
     \\ \\

    \theutterance \stepcounter{utterance}  
    & & \multicolumn{4}{p{0.6\linewidth}}{
        \cellcolor[rgb]{0.9,0.9,0.9}{
            \makecell[{{p{\linewidth}}}]{
                \texttt{\tiny{[P1$\langle$GM]}}
                \texttt{ARGOMENTAZIONE: \{'La proposta include l'oggetto B09, che sembra essere un buon compromesso tra sforzo e importanza. La proposta include oggetti con un buon equilibrio tra sforzo e importanza, e potrebbe essere una buona opzione per entrambi. Tuttavia, potrei cercare di migliorarla ulteriormente, ad esempio sostituendo uno degli oggetti con un altro che abbia un miglior equilibrio tra sforzo e importanza.'\}} \\
\\ 
\texttt{RIFIUTO: \{'B00', 'A08', 'B60', 'B09'\}} \\
\\ 
\texttt{PROPOSTA: \{'A08', 'B60', 'B09', 'A23'\}} \\
            }
        }
    }
    & \\ \\

    \theutterance \stepcounter{utterance}  
    & \multicolumn{4}{p{0.6\linewidth}}{
        \cellcolor[rgb]{0.9,0.9,0.9}{
            \makecell[{{p{\linewidth}}}]{
                \texttt{\tiny{[P1$\rangle$GM]}}
                \texttt{RAGIONAMENTO STRATEGICO: \{'L'altro giocatore ha rifiutato la mia proposta precedente e ha fatto una nuova proposta. Sto analizzando la nuova proposta e sto cercando di capire come rispondere. La proposta include l'oggetto A23, che ha un sforzo di 991 e un'importanza di 262. Questo oggetto ha un sforzo molto alto, quindi devo considerare se è possibile sostituirlo con un altro oggetto che abbia un miglior equilibrio tra sforzo e importanza. Inoltre, la proposta non include più l'oggetto B00, che aveva un'importanza alta.'\}} \\
\\ 
\texttt{ARGOMENTAZIONE: \{'La proposta dell'altro giocatore include l'oggetto A23, che ha un sforzo molto alto. In alternativa, potrei suggerire di sostituirlo con l'oggetto B12, che ha un sforzo di 192 e un'importanza di 915. Questo oggetto potrebbe essere un buon sostituto per l'A23, considerando il miglior equilibrio tra sforzo e importanza. Inoltre, la proposta include già gli oggetti B60 e B09, che sono entrambi buone scelte.'\}} \\
\\ 
\texttt{PROPOSTA: \{'A08', 'B60', 'B09', 'B12'\}} \\
            }
        }
    }
    & & \\ \\

    \theutterance \stepcounter{utterance}  
    & & \multicolumn{4}{p{0.6\linewidth}}{
        \cellcolor[rgb]{0.9,0.9,0.9}{
            \makecell[{{p{\linewidth}}}]{
                \texttt{\tiny{[GM$\rangle$P2]}}
                \texttt{ARGOMENTAZIONE: \{'La proposta dell'altro giocatore include l'oggetto A23, che ha un sforzo molto alto. In alternativa, potrei suggerire di sostituirlo con l'oggetto B12, che ha un sforzo di 192 e un'importanza di 915. Questo oggetto potrebbe essere un buon sostituto per l'A23, considerando il miglior equilibrio tra sforzo e importanza. Inoltre, la proposta include già gli oggetti B60 e B09, che sono entrambi buone scelte.'\}} \\
\\ 
\texttt{PROPOSTA: \{'A08', 'B60', 'B09', 'B12'\}} \\
            }
        }
    }
    & \\ \\

    \theutterance \stepcounter{utterance}  
    & & & \multicolumn{4}{p{0.6\linewidth}}{
        \cellcolor[rgb]{0.9,0.9,0.9}{
            \makecell[{{p{\linewidth}}}]{
                \texttt{\tiny{[GM$\langle$P2]}}
                \texttt{RAGIONAMENTO STRATEGICO: \{'La proposta dell'altro giocatore include l'oggetto B12, che ha un sforzo di 192 e un'importanza di 915, il che sembra essere un ottimo compromesso. La proposta include gli oggetti A08, B60, B09 e B12, con un sforzo totale di 515 + 101 + 501 + 192 = 1309. La proposta sembra essere un buon equilibrio tra sforzo e importanza, e potrebbe essere una buona opzione per entrambi. Considerando la mia strategia, credo che questa proposta sia accettabile e potrebbe essere un buon punto di partenza per raggiungere un accordo.'\}} \\
\\ 
\texttt{ARGOMENTAZIONE: \{'La proposta include l'oggetto B12, che sembra essere un ottimo compromesso tra sforzo e importanza. La proposta include oggetti con un buon equilibrio tra sforzo e importanza, e potrebbe essere una buona opzione per entrambi. Sono disposto ad accettare questa proposta, considerando che rappresenta un buon compromesso tra le nostre esigenze.'\}} \\
\\ 
\texttt{ACCORDO: \{'A08', 'B60', 'B09', 'B12'\}} \\
            }
        }
    }
     \\ \\

    \theutterance \stepcounter{utterance}  
    & & & \multicolumn{2}{p{0.3\linewidth}}{
        \cellcolor[rgb]{0.9,0.9,0.9}{
            \makecell[{{p{\linewidth}}}]{
                \texttt{\tiny{[GM$|$GM]}}
                \texttt{game successful} \\
            }
        }
    }
    & & \\ \\

    \theutterance \stepcounter{utterance}  
    & & & \multicolumn{2}{p{0.3\linewidth}}{
        \cellcolor[rgb]{0.9,0.9,0.9}{
            \makecell[{{p{\linewidth}}}]{
                \texttt{\tiny{[GM$|$GM]}}
                \texttt{end game} \\
            }
        }
    }
    & & \\ \\

\end{supertabular}
}

\end{document}
